\appendix
\clearpage
\phantomsection
\refstepcounter{section}
\addcontentsline{toc}{section}{APÊNDICE A -- EXEMPLOS COMPLETOS DE INTERAÇÃO}
\begin{center}
\textbf{APÊNDICE A -- EXEMPLOS COMPLETOS DE INTERAÇÃO}
\end{center}
\label{sec:apendice-exemplos}

Este apêndice reúne exemplos completos de interação utilizados para complementar a demonstração empírica apresentada na seção de experimentos. O objetivo não é substituir a avaliação sistemática, mas mostrar, em formato legível, como a solução se comporta em um caso de bom desempenho auditável e em um caso que explicita limitações reais de escopo e atribuição.

\subsection{Caso 1: resposta auditável em consulta temática delimitada}\label{sec:anexo-caso-bom}

\textbf{Pergunta}

Quais são os principais argumentos de Paulo Paim contra a capitalização da Previdência e que exemplos ele usa para sustentar essa posição?

\textbf{Resposta}

\textbf{Principais argumentos de Paulo Paim contra a capitalização da Previdência}

\begin{enumerate}
  \item \textbf{``Capitalização'' equivale, na prática, a ``poupança individual'', o que fragiliza a proteção previdenciária.}

  Ele afirma que ``regime de capitalização me preocupa muito'' porque ``significa uma poupança individual'', em que ``quem vai cuidar do meu dinheiro é aquele fundo'' e ``se o fundo aplicar mal, o resultado é zero'' em rendimento.

  \item \textbf{Perde-se a lógica solidária, e sobra apenas uma parcela do trabalhador.}

  Ele argumenta que se retiram diversas fontes de financiamento e se deixa ``só 10\% do empregado'', dizendo que ``claro que não vai dar'', pois ``tu tens 100, tira 90''.

  \item \textbf{Com esse modelo, a aposentadoria tende a ser muito baixa e pode não garantir continuidade para familiares.}

  Ele descreve que, na capitalização, ``se um deles fez o depósito e aos 85 anos o cidadão morreu, sabe qual o vencimento que ela fica? Zero!'', em referência à ausência de pensão. Também sustenta que, com contribuição reduzida, o trabalhador terminaria ``com exatamente um quarto, R\$750''.

  \item \textbf{Exemplos internacionais mostram que o sistema gerou problemas.}

  Ele diz que ``o sistema da capitalização já foi adotado no Chile, na Colômbia, no México e no Peru'' e que ``todos eles estão voltando atrás'', citando aposentadorias muito baixas e cobertura restrita. Como exemplo do Chile, menciona inclusive casos de ``suicídio entre os idosos'' e a ``média da aposentadoria no Chile'' como equivalente a ``meio salário mínimo''.

  \item \textbf{Ele usa dados e cenários para ilustrar o impacto sobre renda e o ganho do sistema financeiro.}

  No debate, cita uma estimativa segundo a qual ``a remuneração dos bancos consumiria R\$105.701,43'', o que equivaleria a mais de 62\% do patrimônio do empregado, deixando o trabalhador com menos. E conclui que ``um sistema sem empregador e com instituição financeira faturando um percentual será um fracasso, é condenar o trabalhador à miséria''.

  \item \textbf{Afirma que não há país que tenha aplicado esse modelo ``com recurso só do trabalhador'' e tido êxito.}

  Ele desafia: ``Se me apresentarem um único país que aplicou a tal de capitalização com o recurso só do trabalhador e que deu certo, eu mudo aqui o meu discurso''.
\end{enumerate}

\textbf{Exemplos usados por Paulo Paim}

\begin{itemize}
  \item Chile: aposentadorias baixas, menção a suicídio entre idosos e retorno parcial atrás do modelo.
  \item Chile, Colômbia, México e Peru: países citados como evidência de insucesso ou revisão do sistema.
  \item Cálculo ilustrativo em reais: estimativa de remuneração bancária consumindo grande parte do patrimônio.
  \item Cenário de pensão: caso em que o cônjuge remanescente ficaria com ``zero''.
  \item Desafio empírico: pedido de um único caso internacional bem-sucedido com recurso apenas do trabalhador.
\end{itemize}

\textbf{Limitações apontadas pela própria resposta}

\begin{itemize}
  \item O contexto não traz uma lista única e consolidada dos argumentos; eles aparecem distribuídos em vários trechos.
  \item A conclusão aparece dispersa em passagens distintas, e não em uma formulação final única.
\end{itemize}

\textbf{Referências retornadas pelo sistema}

\begin{itemize}
  \item Paulo Paim (PT/RS), 21/05/2019, \url{https://www25.senado.leg.br/web/atividade/pronunciamentos/-/p/texto/453686}
  \item Paulo Paim (PT/RS), 20/05/2019, \url{https://www25.senado.leg.br/web/atividade/pronunciamentos/-/p/texto/453766}
  \item Paulo Paim (PT/RS), 08/02/2019, \url{https://www25.senado.leg.br/web/atividade/pronunciamentos/-/p/texto/450736}
  \item Paulo Paim (PT/RS), 19/02/2019, \url{https://www25.senado.leg.br/web/atividade/pronunciamentos/-/p/texto/450992}
  \item Paulo Paim (PT/RS), 07/02/2019, \url{https://www25.senado.leg.br/web/atividade/pronunciamentos/-/p/texto/450672}
\end{itemize}